\documentclass{beamer}
\usetheme{tsubame}

\title{タイトル}
\subtitle{サブタイトル}
\author{著者}
\institute{所属}

\begin{document}

\maketitle

\begin{frame}{概要}
    \begin{myitemlarge}
        \item[背景]\hspace{0pt}\\[-1.3em]
        myitemlarge 環境は概要・まとめのスライド用
        \item[提案]\hspace{0pt}\\[-1.3em]
        見出しを大きく，項目の間を広く取る
        \item[結果]\hspace{0pt}\\[-1.3em]
        3項目くらいまでが収まりよい
    \end{myitemlarge}
\end{frame}

\section{基本の要素}

\begin{frame}{箇条書きと強調}
    \begin{itemize}
        \item 本文は Science Blue
        \item \alert{\textbackslash alert} で強調するとアクセントのオレンジ
        \item 数式は New Computer Modern Math
    \end{itemize}
    \[
        f(x) = \sum_{i=1}^{n} w_i \phi_i(x)
    \]
\end{frame}

\begin{frame}{description 環境}
    \subsectionhead{myitem}
    \begin{myitem}
        \item[見出し1]\hspace{0pt}\\[-1.3em]
        詳細スライド用．見出しの下に内容を書く
        \item[見出し2]\hspace{0pt}\\[-1.3em]
        内容2
    \end{myitem}
\end{frame}

\begin{frame}{ブロック}
    \begin{block}{定義}
        ブロックの地はエクステンディッドカラーのバックグラウンド2色
    \end{block}
    \begin{alertblock}{注意}
        alertblock は見出しがオレンジになる
    \end{alertblock}
\end{frame}

\section{図}

\begin{frame}{tsubamebox と tsubamearrow}
    \begin{center}
        \begin{tikzpicture}
            \node[tsubamebox] (a) {入力};
            \node[tsubamebox, right=2cm of a] (b) {出力};
            \draw[tsubamearrow] (a) -- (b);
        \end{tikzpicture}
    \end{center}
\end{frame}

\begin{frame}[standout]
    ご清聴ありがとうございました
\end{frame}

\end{document}
